% GENERATED FILE. Edit canonical lessons, then run npm run generate:books.
% canonical-source-hash: fnv1a64:5815ed67ae1e55e7
% canonical-lessons: FR-C482-eternuer, FR-C482-allergie, FR-C482-pollen, FR-C482-demangeaison, FR-C482-gratter

\chapter{Éternuements et allergies}
\label{ch:fr-eternuements-et-allergies}
\hlchaptermodality{482}

\begin{chapteropening}
\textbf{By the end of this chapter:} \emph{I can talk about allergies, sneezing and itching.}

Complete the last lesson of chapter 482: I can talk about allergies, sneezing and itching.
\end{chapteropening}

\section[éternuer]{éternuer — to sneeze}

\label{lesson:FR-C482-eternuer}



\begin{quote}
Before the new one: say the French for splint, then the French for physiotherapist.
\end{quote}

\subsection*{You'll want to know: éternuer}
\textbf{éternuer} — \textquotedblleft{}to sneeze\textquotedblright{}.

\begin{quote}
Au printemps, j'éternue tout le temps.
\end{quote}

\begin{quote}
\emph{In spring I sneeze all the time.}
\end{quote}

\emph{À vos souhaits !} — bless you! (said after a sneeze).

\begin{cousinweb}
\emph{éternuer} is inherited, through Old French \emph{esternuer}, from Latin \emph{sternūtāre} \textquotedblleft{}to sneeze\textquotedblright{} (EtymDB), built on \emph{sternuere} (Lewis \& Short).
\end{cousinweb}

\begin{grammarlens}[title={What you've built}]
One more word for telling, explaining and describing at B1.
\end{grammarlens}

\subsection*{Your turn}
\begin{itemize}
\raggedright
  \item \emph{Say it:} \emph{éternuer}
  \item \emph{Say it:} \emph{éternuer}, once more
  \item \emph{Say it:} a sentence of your own with \emph{éternuer}
  \item \emph{Recall:} say the French for splint, then the French for physiotherapist, then say \emph{éternuer} again
\end{itemize}

\begin{tcolorbox}[breakable,colback=teal!4,colframe=teal!35!black,title={Before you move on}]
What does éternuer mean? (\textquotedblleft{}To sneeze\textquotedblright{}.) Say it once more.
\end{tcolorbox}
\section[l'allergie]{l'allergie — allergy}

\label{lesson:FR-C482-allergie}



\begin{quote}
Before the new one: say the French for physiotherapist, then the French for to sneeze.
\end{quote}

\subsection*{You'll want to know: l'allergie}
\textbf{l'allergie} — \textquotedblleft{}allergy\textquotedblright{}.

\begin{quote}
Mon fils a une allergie aux chats.
\end{quote}

\begin{quote}
\emph{My son has an allergy to cats.}
\end{quote}

Feminine. \emph{Être allergique à} — to be allergic to.

\begin{cousinweb}
\emph{l'allergie} was borrowed from German \emph{Allergie}, a modern learned word built from Greek \emph{állos} \textquotedblleft{}other\textquotedblright{} and \emph{érgon} \textquotedblleft{}work, action\textquotedblright{}.
\end{cousinweb}

\begin{grammarlens}[title={What you've built}]
One more word for telling, explaining and describing at B1.
\end{grammarlens}

\subsection*{Your turn}
\begin{itemize}
\raggedright
  \item \emph{Say it:} \emph{l'allergie}
  \item \emph{Say it:} \emph{l'allergie}, once more
  \item \emph{Say it:} a sentence of your own with \emph{l'allergie}
  \item \emph{Recall:} say the French for physiotherapist, then the French for to sneeze, then say \emph{l'allergie} again
\end{itemize}

\begin{tcolorbox}[breakable,colback=teal!4,colframe=teal!35!black,title={Before you move on}]
What does l'allergie mean? (\textquotedblleft{}Allergy\textquotedblright{}.) Say it once more.
\end{tcolorbox}
\section[le pollen]{le pollen — pollen}

\label{lesson:FR-C482-pollen}



\begin{quote}
Before the new one: say the French for to sneeze, then the French for allergy.
\end{quote}

\subsection*{You'll want to know: le pollen}
\textbf{le pollen} — \textquotedblleft{}pollen\textquotedblright{}.

\begin{quote}
En mai, il y a beaucoup de pollen dans l'air.
\end{quote}

\begin{quote}
\emph{In May there is a lot of pollen in the air.}
\end{quote}

The final \emph{n} is sounded: \emph{pol-lène}.

\begin{cousinweb}
\emph{le pollen} is a learned borrowing from Latin \emph{pollen} \textquotedblleft{}fine flour, powder\textquotedblright{} (Lewis \& Short).
\end{cousinweb}

\begin{grammarlens}[title={What you've built}]
One more word for telling, explaining and describing at B1.
\end{grammarlens}

\subsection*{Your turn}
\begin{itemize}
\raggedright
  \item \emph{Say it:} \emph{le pollen}
  \item \emph{Say it:} \emph{le pollen}, once more
  \item \emph{Say it:} a sentence of your own with \emph{le pollen}
  \item \emph{Recall:} say the French for to sneeze, then the French for allergy, then say \emph{le pollen} again
\end{itemize}

\begin{tcolorbox}[breakable,colback=teal!4,colframe=teal!35!black,title={Before you move on}]
What does le pollen mean? (\textquotedblleft{}Pollen\textquotedblright{}.) Say it once more.
\end{tcolorbox}
\section[la démangeaison]{la démangeaison — itch, itching}

\label{lesson:FR-C482-demangeaison}



\begin{quote}
Before the new one: say the French for allergy, then the French for pollen.
\end{quote}

\subsection*{You'll want to know: la démangeaison}
\textbf{la démangeaison} — \textquotedblleft{}itch, itching\textquotedblright{}.

\begin{quote}
Cette crème calme les démangeaisons.
\end{quote}

\begin{quote}
\emph{This cream soothes the itching.}
\end{quote}

The verb is \emph{démanger}: \emph{ça me démange} — it itches.

\begin{cousinweb}
\emph{la démangeaison} is built on \emph{démanger} \textquotedblleft{}to itch\textquotedblright{} with \emph{-aison} (EtymDB); \emph{démanger} is \emph{manger} \textquotedblleft{}to eat\textquotedblright{}, from Latin \emph{mandūcāre} \textquotedblleft{}to chew\textquotedblright{}, with \emph{dé-}: the skin feels eaten.
\end{cousinweb}

\begin{grammarlens}[title={What you've built}]
One more word for telling, explaining and describing at B1.
\end{grammarlens}

\subsection*{Your turn}
\begin{itemize}
\raggedright
  \item \emph{Say it:} \emph{la démangeaison}
  \item \emph{Say it:} \emph{la démangeaison}, once more
  \item \emph{Say it:} a sentence of your own with \emph{la démangeaison}
  \item \emph{Recall:} say the French for allergy, then the French for pollen, then say \emph{la démangeaison} again
\end{itemize}

\begin{tcolorbox}[breakable,colback=teal!4,colframe=teal!35!black,title={Before you move on}]
What does la démangeaison mean? (\textquotedblleft{}Itch, itching\textquotedblright{}.) Say it once more.
\end{tcolorbox}
\section[gratter]{gratter — to scratch; to itch}

\label{lesson:FR-C482-gratter}



\begin{quote}
Before the new one: say the French for pollen, then the French for itch.
\end{quote}

\subsection*{You'll want to know: gratter}
\textbf{gratter} — \textquotedblleft{}to scratch; to itch\textquotedblright{}.

\begin{quote}
Le chien se gratte derrière l'oreille.
\end{quote}

\begin{quote}
\emph{The dog is scratching behind its ear.}
\end{quote}

\emph{Se gratter} — to scratch oneself. \emph{Ce pull gratte} — this jumper is itchy.

\begin{cousinweb}
\emph{gratter} comes from Frankish \emph{krattōn} \textquotedblleft{}to scratch\textquotedblright{}, a Germanic word related to German \emph{kratzen} (EtymDB).
\end{cousinweb}

\begin{grammarlens}[title={What you've built}]
One more word for telling, explaining and describing at B1.
\end{grammarlens}

\subsection*{Your turn}
\begin{itemize}
\raggedright
  \item \emph{Say it:} \emph{gratter}
  \item \emph{Say it:} \emph{gratter}, once more
  \item \emph{Say it:} a sentence of your own with \emph{gratter}
  \item \emph{Recall:} say the French for pollen, then the French for itch, then say \emph{gratter} again
\end{itemize}

\begin{tcolorbox}[breakable,colback=teal!4,colframe=teal!35!black,title={Before you move on}]
What does gratter mean? (\textquotedblleft{}To scratch; to itch\textquotedblright{}.) Say it once more.
\end{tcolorbox}
